A comunicação é uma das capacidades mais elementares do ser humano. É por meio dela que uma pessoa expressa necessidades imediatas, como dor, sede ou desconforto, participa das decisões sobre o próprio cuidado e mantém seus vínculos familiares e sociais. Diversas condições neurológicas, no entanto, comprometem de forma progressiva ou súbita os movimentos necessários para falar, escrever ou gesticular, ao mesmo tempo em que preservam a consciência e a capacidade cognitiva. Nesses casos, a pessoa continua tendo o que dizer, mas perde os meios de fazê-lo.

A \gls{ELA} é um dos exemplos mais representativos desse cenário. Trata-se de uma doença neurodegenerativa que leva à paralisia gradual da musculatura voluntária, incluindo a responsável pela fala, mas que costuma poupar os movimentos oculares até fases avançadas (Seção~\ref{sec:ela}). O olhar torna-se, assim, o principal canal de comunicação disponível ao paciente. Tecnologias de rastreamento ocular (\textit{eye-tracking}) exploram esse canal ao estimar para qual ponto da tela a pessoa está olhando e converter esse olhar em comandos, o que permite selecionar letras, palavras e frases sem nenhum movimento das mãos. As soluções comerciais que oferecem esse recurso, porém, dependem de câmeras infravermelhas e de equipamentos dedicados, com custo elevado para a realidade brasileira (Seção~\ref{sec:solucoes_existentes}).

Este trabalho propõe um aplicativo para \textit{tablets} Android comuns que realiza o rastreamento ocular apenas com a câmera frontal do aparelho e oferece uma interface de digitação controlada pelo olhar, com sugestão de palavras e síntese de voz, sem depender de \textit{hardware} adicional nem de conexão com a internet. As seções a seguir apresentam a doença que motiva o trabalho, os problemas enfrentados por quem perde a capacidade de se comunicar, as soluções existentes e suas limitações, os objetivos propostos e a organização deste documento.

\section{Esclerose Lateral Amiotrófica} \label{sec:ela}

A \gls{ELA} é uma doença neurodegenerativa progressiva caracterizada pela degeneração dos neurônios motores superiores, que partem do córtex cerebral em direção ao tronco encefálico e à medula espinhal, e dos neurônios motores inferiores, que partem do tronco encefálico e da medula em direção aos músculos~\cite{hardiman2017als}. A perda desses neurônios leva à fraqueza e, por fim, à paralisia da musculatura voluntária. Na grande maioria dos casos, acima de 90\%, não há histórico familiar da doença, e o processo degenerativo começa bem antes de se tornar perceptível: estima-se que, no momento em que aparece o primeiro sintoma, mais de 80\% dos neurônios motores já tenham se perdido~\cite{pcdt_ela2020}.

A forma de início varia entre os pacientes. No início espinhal, os primeiros sinais são a fraqueza e a atrofia dos músculos dos membros; no início bulbar, a doença se manifesta primeiro na musculatura da boca, da língua e da garganta, com disartria (dificuldade de articular a fala) e disfagia (dificuldade de deglutir)~\cite{hardiman2017als}. No levantamento nacional conduzido por Dietrich-Neto \textit{et al.}~\cite{dietrichneto2000als}, o início foi espinhal em 69\% dos pacientes, bulbar em 18,5\% e simultâneo nas duas regiões em 11,7\%. Seja qual for o ponto de partida, a degeneração avança para outras regiões do sistema motor, que vão da bulbar à lombar, passando pelas regiões cervical e torácica~\cite{pcdt_ela2020}. A disartria é o sintoma inicial em cerca de 30\% dos casos e está presente em mais de 80\% dos pacientes ao longo da doença, podendo evoluir até a anartria, a perda completa da capacidade de articular a fala~\cite{hardiman2017als}.

A sobrevida média é de 3 a 5 anos, e a morte decorre, na maioria dos casos, de insuficiência respiratória, causada pelo enfraquecimento dos músculos da respiração~\cite{pcdt_ela2020}. Apenas 8\% a 16\% dos pacientes chegam a dez anos de doença sem ventilação mecânica prolongada, ao passo que o suporte ventilatório pode estender a sobrevida para 15 anos ou mais~\cite{pcdt_ela2020}. Isso significa que o suporte respiratório prolonga justamente as fases mais avançadas da doença, nas quais a fala e o movimento dos membros já foram perdidos.

No Brasil, os dados epidemiológicos sobre a \gls{ELA} são escassos, e a maior parte dos estudos tem alcance estadual ou municipal~\cite{barbalho2023als}. O primeiro levantamento nacional, realizado em 1998 a partir de formulários enviados a 2.505 neurologistas, reuniu dados de 443 pacientes: 58,5\% eram homens, a idade média de início dos sintomas foi de 52 anos e a incidência ajustada por idade atingiu o pico na faixa de 65 a 74 anos~\cite{dietrichneto2000als}. Um estudo posterior, baseado em 8.942 atestados de óbito emitidos entre 2004 e 2013 em todo o país, encontrou idade média ao óbito de 62,7 anos, predominância masculina (razão de 1,3:1) e taxa de mortalidade ajustada entre 0,61 e 0,89 por 100 mil pessoas-ano na população acima de 20 anos, chegando a valores entre 1,77 e 2,3 por 100 mil pessoas-ano acima de 45 anos~\cite{moura2016als}. Em Porto Alegre, a prevalência estimada foi de 5,0 casos por 100 mil habitantes em 2010~\cite{linden2013als}.

Dois aspectos da doença são centrais para este trabalho. O primeiro é a preservação da sensibilidade e, em muitos casos, da cognição. Alterações sensitivas não fazem parte do quadro e, salvo quando explicadas por outra doença associada, afastam o diagnóstico de \gls{ELA}. Já no campo cognitivo, avaliações detalhadas podem identificar alterações em até metade dos pacientes, mas a demência é pouco frequente~\cite{pcdt_ela2020}: a variante comportamental da demência frontotemporal acompanha cerca de 13\% dos casos~\cite{hardiman2017als}. O segundo aspecto é a relativa resistência dos neurônios que controlam os movimentos dos olhos. Os neurônios motores oculares, que comandam os músculos dos olhos, estão entre os subgrupos relativamente resistentes à degeneração e persistem ao longo da doença; por isso, dispositivos de rastreamento ocular são usados para que pacientes paralisados, que já não conseguem falar, possam se comunicar~\cite{nijssen2017motor}. Essa preservação não é absoluta: alterações oculomotoras já foram descritas na \gls{ELA}, sobretudo em pacientes com sobrevida prolongada~\cite{sharma2011oculomotor}. Nas fases avançadas, as técnicas de \gls{CAA} são o tratamento de escolha para a perda da fala e podem melhorar a qualidade de vida do paciente~\cite{hardiman2017als}, e é nesse contexto que o rastreamento ocular ganha importância.

A \gls{ELA} não é a única condição em que o olhar permanece como principal canal de comunicação. Na síndrome do encarceramento (\textit{locked-in syndrome}), o paciente está acordado e consciente, mas não consegue falar nem mover os membros ou a face; a causa mais comum são lesões agudas na porção ventral da ponte, no tronco encefálico, e a comunicação passa a depender de códigos baseados nos movimentos dos olhos. Uma vez estabilizado clinicamente, o paciente pode usar tecnologias controladas pelo olhar para escrever em um processador de texto acoplado a um sintetizador de voz e acessar a internet~\cite{laureys2005locked}. O controle pelo olhar também vem sendo investigado como forma de acesso à comunicação para pessoas com paralisia cerebral e comprometimento motor grave, ainda que as evidências sejam limitadas~\cite{karlsson2018eyegaze}, e famílias de pessoas com síndrome de Rett relatam o uso de sistemas de rastreamento ocular como meio de \gls{CAA}~\cite{townend2016rett}. Embora a \gls{ELA} motive este trabalho, a solução proposta pode beneficiar qualquer pessoa que preserve o controle dos movimentos oculares e tenha perdido os demais meios de comunicação.

\section{Problemas Enfrentados} \label{sec:problemas}

A perda da fala é uma das consequências mais frequentes da \gls{ELA}. Estima-se que de 80\% a 95\% dos pacientes apresentem alterações da fala ao longo da doença~\cite{pcdt_ela2020} e que, em algum momento, deixem de conseguir suprir suas necessidades diárias de comunicação apenas com a fala natural~\cite{beukelman2011communication}. Essa perda costuma ser rápida: quando a inteligibilidade da fala começa a cair, ela tende a se deteriorar em pouco tempo, o que deixa uma janela curta para a adoção de meios alternativos de comunicação~\cite{beukelman2011communication}. Nas fases avançadas, a fraqueza também atinge as mãos e os braços, e muitos pacientes chegam à tetraplegia associada à anartria~\cite{caligari2013eye}: escrever à mão ou digitar em um teclado convencional deixam, então, de ser alternativas. A necessidade de se comunicar, no entanto, permanece em todas as esferas da vida: no autocuidado, nas consultas e tratamentos, no convívio social, na vida comunitária e, para alguns, no trabalho~\cite{beukelman2011communication}.

Com a progressão da doença, o paciente perde a autonomia em atividades cotidianas. Não por acaso, a principal escala funcional usada para acompanhar a progressão da \gls{ELA} avalia justamente tarefas como escrever à mão, cortar os alimentos, vestir-se e cuidar da higiene e virar-se na cama~\cite{cedarbaum1999alsfrs}. O cuidado diário recai principalmente sobre a família. Em um estudo com oito cuidadores de pessoas com \gls{ELA} em Recife, a maioria apresentou sobrecarga moderada e quase todos moravam com o paciente~\cite{almeida2017sobrecarga}. A família também costuma assumir o suporte à comunicação: em um levantamento citado por Beukelman \textit{et al.}~\cite{beukelman2011communication}, 96\% das pessoas que davam apoio diário ao uso dos recursos de comunicação de pacientes com \gls{ELA} eram familiares, em sua maioria sem formação técnica. Conforme a fala perde inteligibilidade, compreender o paciente exige cada vez mais esforço de quem o escuta, e a dificuldade de se fazer entender tende a levá-lo ao isolamento~\cite{pcdt_ela2020}. Sem um meio eficiente de se expressar, necessidades simples, como relatar dor, pedir água ou pedir para ser reposicionado, dependem de que o cuidador as adivinhe ou as interprete.

O protocolo do Ministério da Saúde aponta a perda da comunicação efetiva como um dos aspectos mais graves da doença, com reflexos psicológicos que alcançam as esferas emocional, social, familiar e profissional, pioram a qualidade de vida e podem até reduzir a sobrevida~\cite{pcdt_ela2020}. A relação entre comunicação e qualidade de vida aparece também de outra forma. Muitos pacientes com \gls{ELA} avaliam sua qualidade de vida global como boa mesmo com o avanço da incapacidade física, o que se explica em parte por uma reavaliação do que consideram importante: as atividades físicas perdem peso e as relações sociais ganham papel central~\cite{hardiman2017als}. Justamente essas relações dependem da capacidade de se comunicar. Por outro lado, a perda de autonomia, a dependência de terceiros e a perda de dignidade aparecem entre os motivos relatados por pacientes que optaram pela morte assistida em locais onde ela é permitida~\cite{hardiman2017als}. Há evidências de que a introdução precoce de dispositivos de comunicação assistiva melhora o bem-estar psicológico e existencial dos pacientes e a percepção de apoio de seus cuidadores~\cite{londral2015quality}.

Sem acesso a recursos tecnológicos, a comunicação se apoia em métodos de baixa tecnologia. Enquanto ainda resta alguma fala, o paciente pode apontar em uma prancha alfabética a primeira letra de cada palavra para ajudar o interlocutor a compreendê-lo~\cite{beukelman2011communication}. Em fases mais avançadas, são comuns códigos de ``sim'' e ``não'' feitos com os olhos, como os empregados na síndrome do encarceramento~\cite{laureys2005locked}, e as pranchas de transferência do olhar (\textit{e-tran}): placas transparentes com grupos de letras, seguradas pelo interlocutor, que observa para onde o paciente olha e decodifica cada letra a partir de dois movimentos oculares~\cite{zhang2017gazespeak}. Esses métodos são baratos, mas lentos: escolher uma única letra leva de 8 a 12 segundos~\cite{zhang2017gazespeak}. Além disso, impõem uma carga cognitiva alta ao interlocutor, que precisa decodificar e memorizar as letras e prever as palavras~\cite{zhang2017gazespeak}. Em consequência, o paciente só consegue se comunicar quando há alguém treinado disponível e disposto a fazer esse papel. Não surpreende que, em um estudo com pessoas com \gls{ELA}, cuidadores e ouvintes não familiarizados, os três grupos tenham demonstrado forte preferência por recursos de \gls{CAA} de alta tecnologia em relação tanto à fala de difícil compreensão quanto a cadernos de comunicação de baixa tecnologia~\cite{beukelman2011communication}.

Os métodos de baixa tecnologia continuam importantes quando não há outra opção disponível~\cite{zhang2017gazespeak}, mas não oferecem ao paciente uma comunicação autônoma e com velocidade compatível com uma conversa. Entre os métodos de acesso a recursos de alta tecnologia, o rastreamento ocular é frequentemente relatado por pessoas com \gls{ELA} como o menos cansativo, e em um dos estudos revisados por Beukelman \textit{et al.}~\cite{beukelman2011communication} ele foi escolhido por 53\% dos participantes porque o movimento dos olhos era o único meio de acesso viável. Em um estudo com 35 usuários de dispositivos de comunicação por rastreamento ocular em fase avançada da doença, com tetraplegia e anartria, esses dispositivos ampliaram a capacidade de comunicação além do obtido com a prancha \textit{e-tran}, com grande impacto positivo na qualidade de vida e alta satisfação dos usuários~\cite{caligari2013eye}. Em outro estudo, com 30 pacientes em fase avançada, o uso diário mediano foi de 300 minutos, principalmente para conversar com familiares e cuidadores, navegar na internet e trocar mensagens; entre os que usavam pouco o dispositivo, as causas mais citadas foram o cansaço ocular e alterações oculomotoras~\cite{spataro2014eye}. A próxima seção apresenta as soluções de rastreamento ocular disponíveis no mercado e as barreiras que limitam o acesso a elas.

\section{Soluções Existentes} \label{sec:solucoes_existentes}

As soluções voltadas a pessoas que perderam a fala se inserem na área da \gls{CAA}, tradução do termo inglês \textit{augmentative and alternative communication}, também chamada de Comunicação Suplementar e Alternativa no protocolo do Ministério da Saúde para a \gls{ELA}~\cite{pcdt_ela2020}. Segundo a \textit{American Speech-Language-Hearing Association} (ASHA), trata-se de uma área de prática clínica voltada a suprir ou compensar limitações na produção ou na compreensão da linguagem, tanto falada quanto escrita~\cite{asha_aac}. A \gls{CAA} é dita aumentativa quando complementa uma fala ainda existente e alternativa quando substitui uma fala ausente ou não funcional~\cite{asha_aac}. Seus recursos vão desde gestos e pranchas de figuras, como os métodos de baixa tecnologia descritos na Seção~\ref{sec:problemas}, até dispositivos geradores de fala: equipamentos eletrônicos que permitem compor mensagens e reproduzi-las em voz sintetizada. Para quem não consegue usar as mãos, esses dispositivos podem ser controlados por métodos de acesso alternativos, como acionadores, movimentos da cabeça e rastreamento ocular~\cite{beukelman2011communication}.

No mercado, os dispositivos de comunicação controlados pelo olhar mais conhecidos são os da Smartbox e os da Tobii Dynavox. A Smartbox oferece a linha de \textit{tablets} Grid Pad, que pode ser equipada com a câmera de rastreamento ocular Lumin-i. A câmera usa emissores infravermelhos que permitem o uso ao ar livre, tolera movimentos da cabeça dentro de um volume de rastreamento de $35 \times 30 \times 65$~cm e tem acurácia declarada de $0{,}5^\circ$~\cite{smartbox_lumini}. A comunicação é feita pelo \textit{software} Grid~3, que oferece comunicação por símbolos e por texto, com teclado virtual e predição de palavras, além de controle do ambiente e acesso ao computador~\cite{grid3}. A Tobii Dynavox oferece a linha TD I-Series, com os modelos TD I-13 e TD I-16: dispositivos dedicados à \gls{CAA}, com rastreador ocular embutido capaz de funcionar inclusive sob luz do sol, uma tela traseira (\textit{partner window}) que exibe o texto ao interlocutor durante a conversa e os \textit{softwares} TD Snap e Communicator~5 pré-instalados~\cite{tobii_iseries}.

Esses equipamentos reúnem, em um aparelho dedicado, rastreamento ocular preciso e capaz de funcionar ao ar livre, comunicação por símbolos e por texto com predição de palavras e voz sintetizada; no caso da TD I-Series, o fabricante destaca ainda a durabilidade e a resistência à água~\cite{tobii_iseries}. O custo, porém, é elevado. Nos Estados Unidos, o Grid Pad 13 custa US\$~8.495 e a câmera Lumin-i, US\$~7.990, o que totaliza US\$~16.485 para o conjunto~\cite{smartbox_precos}; o preço sugerido do TD I-13 é de US\$~15.565~\cite{tobii_preco}. Esses valores não incluem frete nem tributos de importação, o que coloca tais dispositivos fora do alcance da maior parte das famílias brasileiras. Além do custo, os rastreadores comerciais exigem \textit{hardware} dedicado, com câmeras e emissores infravermelhos específicos, e um suporte que mantenha o equipamento relativamente fixo em relação ao usuário~\cite{zhang2017gazespeak}.

Ao mesmo tempo, \textit{tablets} e celulares comuns já trazem uma câmera frontal e capacidade de processamento suficiente para executar modelos de visão computacional no próprio aparelho. Bibliotecas como o ML Kit, do Google, detectam pontos do rosto em tempo real diretamente no dispositivo móvel~\cite{mlkit_facemesh}, e trabalhos acadêmicos mostraram que é possível estimar a direção do olhar a partir de imagens de câmeras comuns, sem iluminação infravermelha, com redes neurais treinadas em grandes conjuntos de imagens. O MPIIGaze, por exemplo, reúne 213.659 imagens de 15 participantes capturadas pela \textit{webcam} durante o uso cotidiano de seus \textit{notebooks}~\cite{zhang2015mpiigaze}, e o iTracker, treinado com dados de mais de 1.450 pessoas, estima o ponto observado na tela em tempo real em um dispositivo móvel, com erro de 2,12~cm em \textit{tablets} após calibração~\cite{krafka2016eye}. Propostas como o GazeSpeak exploram esse caminho para a comunicação de pessoas com \gls{ELA}, ao interpretar gestos oculares capturados pela câmera de um celular~\cite{zhang2017gazespeak}.

A estimativa do olhar com câmeras comuns ainda não alcança a precisão dos rastreadores infravermelhos dedicados: enquanto o iTracker erra cerca de 2~cm em \textit{tablets}, a acurácia de $0{,}5^\circ$ declarada para o Lumin-i corresponde a cerca de meio centímetro na tela, à distância típica de uso. Em compensação, exige apenas um aparelho de consumo comum, e não um equipamento especializado. Além disso, o erro tem peso menor em telas maiores, como as dos \textit{tablets}, que comportam alvos maiores e mais espaçados. É nessa lacuna que se insere este trabalho.

\section{Objetivos} \label{sec:objetivos}

O objetivo geral deste trabalho é desenvolver um aplicativo de comunicação controlado pelo olhar que funcione em \textit{tablets} Android comuns, utilizando apenas a câmera frontal do aparelho, sem \textit{hardware} adicional e sem depender de conexão com a internet, de modo a oferecer uma alternativa de baixo custo aos dispositivos comerciais descritos na Seção~\ref{sec:solucoes_existentes}. O trabalho parte da construção de uma versão mínima funcional do sistema, que é então avaliada e aprimorada a partir dos resultados obtidos.

Para alcançar esse objetivo, foram definidos os seguintes objetivos específicos:

\begin{enumerate}
    \item Implementar o rastreamento ocular a partir das imagens da câmera frontal, combinando a detecção de pontos do rosto, a localização da pupila por uma rede neural e uma etapa de calibração por usuário que relacione a posição dos olhos ao ponto observado na tela;
    \item Desenvolver uma interface de digitação que possa ser operada exclusivamente pelo olhar, convertendo a direção do olhar em comandos de navegação e seleção;
    \item Reduzir o esforço de digitação por meio da sugestão de palavras, utilizando um modelo de linguagem baseado em n-gramas construído a partir de um conjunto de frases típicas de comunicação;
    \item Converter o texto produzido em fala com um modelo de síntese de voz em português executado no próprio aparelho, sem envio de dados a servidores externos;
    \item Registrar métricas de desempenho da digitação e avaliar o sistema em testes com usuários, analisando a velocidade de comunicação, o esforço necessário e a contribuição das sugestões de palavras.
\end{enumerate}

\section{Estrutura do Trabalho} \label{sec:estrutura}

Este documento está organizado em cinco capítulos, incluindo esta introdução.

O \refCap{2_Referencial_Teorico} apresenta o referencial teórico que fundamenta o trabalho. Inicialmente, descreve as tecnologias de desenvolvimento utilizadas: o \textit{framework} Flutter, a linguagem Dart, o sistema operacional Android e o acesso a recursos nativos do aparelho, como a câmera. Em seguida, aborda os conceitos de inteligência artificial e aprendizado de máquina empregados no rastreamento ocular, incluindo o detector de objetos YOLO e a detecção de pontos do rosto pelo ML Kit, além do filtro de Kalman, dos modelos de linguagem utilizados no preenchimento automático de palavras, das métricas de avaliação de entrada de texto e da síntese de voz.

O \refCap{3_Metodologia_e_Desenvolvimento} descreve a metodologia e o desenvolvimento do aplicativo: o ambiente e as ferramentas de desenvolvimento, a arquitetura do \textit{software} e seus módulos, a interface visual e a implementação de cada componente do sistema, a saber, o rastreamento ocular, com os conjuntos de dados, os modelos, a calibração e a filtragem do sinal, o preenchimento automático de palavras, a síntese de voz e o registro das métricas de digitação. O capítulo também relata como ferramentas de inteligência artificial foram utilizadas no processo de desenvolvimento.

O \refCap{4_Resultados} apresenta os resultados obtidos, desde os testes realizados durante o desenvolvimento até a avaliação do sistema com usuários, e analisa as métricas de digitação coletadas.

Por fim, o \refCap{5_Conclusoes} retoma os objetivos propostos, discute em que medida foram alcançados e aponta possibilidades de trabalhos futuros.
